% Authored: release process.

\chapter{Releases}
\label{ch:releases}

\section{Versioning}

The project version at this revision is \texttt{\projectversion}. Version
numbers follow a \texttt{MAJOR.MINOR.PATCH} form and the changelog history
shows minor increments carrying new features, which is consistent with
semantic versioning~\cite{semver}. The repository states no versioning policy
document, so the scheme is inferred from the artifacts rather than quoted from
a rule.

\noevidence

The version is declared in four places that must agree:

\begin{longtable}{p{0.34\textwidth} p{0.30\textwidth} p{0.28\textwidth}}
\toprule
File & Field & Purpose \\
\midrule
\endfirsthead
\toprule
File & Field & Purpose \\
\midrule
\endhead
\bottomrule
\endlastfoot
\texttt{pyproject.toml} & \texttt{[tool.poetry] version} & Distribution metadata \\
\texttt{src/.../\_\_init\_\_.py} & \texttt{\_\_version\_\_} & Reported by the API and CLI \\
\texttt{CITATION.cff} & \texttt{version}, \texttt{date-released} & Software citation \\
\texttt{.zenodo.json} & \texttt{version} & Archive deposit metadata \\
\end{longtable}

Nothing enforces the agreement automatically; the release workflow validates
that the two metadata files \emph{parse}, not that their versions match
\srcfile{pyproject.toml}. See \cref{ch:improvements}.

\section{Changelog}

\srcfile{CHANGELOG.md} is maintained by hand in a Keep a
Changelog~\cite{keepachangelog} style: an \texttt{Unreleased} section at the
top, then one section per version with a release date and \texttt{Added},
\texttt{Changed}, and \texttt{Fixed} subsections. Entries are written in terms
of user-visible capability rather than commits.

Separately, \srcfile{.github/release.yml} configures GitHub's automatic release
notes: six categories (Features, Fixes, Documentation, Dependencies,
Maintenance, Other), with \texttt{dependabot[bot]} and anything labelled
\texttt{ignore-for-release} excluded. The two mechanisms coexist --- the
hand-written changelog is the narrative, the generated notes are the pull
request digest.

\section{Release workflow}

Triggered by pushing a tag matching \texttt{v*} or by manual dispatch, on
Python 3.12 and Node 24:

\begin{enumerate}[nosep]
  \item check out the repository;
  \item install Poetry (pinned) and the Python dev dependencies;
  \item install frontend dependencies with \texttt{npm ci};
  \item validate \srcfile{.zenodo.json} parses as JSON;
  \item validate \srcfile{CITATION.cff} parses as YAML;
  \item run the deterministic demo;
  \item build the Python distribution;
  \item build the frontend;
  \item upload \srcfile{dist/*} as the \texttt{python-dist} artifact;
  \item upload \srcfile{frontend/dist} as the \texttt{frontend-dist} artifact.
\end{enumerate}

The workflow requests \texttt{contents: read} and references no secret. It
therefore \emph{validates and builds} a release; it does not publish one. There
is no PyPI upload, no container registry push, no GitHub Release creation, and
no deployment step anywhere in \srcfile{.github/workflows/}. Distribution
beyond the workflow artifacts is a manual step, and the repository does not
document it.

\noevidence

\section{Archival and citation}

The project is archived on Zenodo with the concept DOI
\texttt{\projectdoi}. \srcfile{CITATION.cff} follows the Citation File
Format~\cite{citationcff} and carries author, ORCID, affiliation, licence,
version, release date, and repository URL; \srcfile{.zenodo.json} carries the
deposit metadata including the related-identifier link back to the repository.
Both are validated on every release run, which is what keeps the archive
deposit from failing on malformed metadata.

\section{Licence}

The project is licensed under \projectlicense, declared in \srcfile{LICENSE}
and repeated in \srcfile{CITATION.cff} and \srcfile{.zenodo.json}.
\srcfile{pyproject.toml} does not declare a \texttt{license} field, so the
built distribution's metadata does not carry it; the licence file itself is the
authoritative statement.
